\section{The actual classical Hall action and its normalization}
\label{classic:sec:action}

We use the trace pairing \((a,b)=-\operatorname{Tr}(ab)/2\) on
\(\mathfrak{su}(2)\), so that unit imaginary quaternions have norm one,
\(\operatorname{Vol}SU(2)=2\pi^2\), and the metric effective group has
\(\operatorname{Vol}SO(3)=\pi^2\).
The exponent convention is \(e^{2\pi i r\,\mathrm{CS}}\), where \(r=k+2\).
On the flat actions \(0,\frac12\), replacing \(k\) by \(r\) leaves the
phase unchanged; in the leading normal integrals the level shift changes
neither the scale nor the leading coefficient. All classical integrations
below are finite residual integrals: the nondegenerate quadratic complement
is integrated invariantly, the finite zero modes are retained, and no
identification with a nonperturbative functional integral is assumed.

Let \(S\) be the genuinely noncentral fibres in an all-negative-centre
Hall chart, \(T=N(S)\), and let \(I\) be the induced forest on the other
vertices. Here the Hall hypothesis on the chart is that \(I\) is balanced
with a perfect matching (and \(S\) has the prescribed noncentral data).
The case \(S=T=\varnothing\) with a perfectly matched tree is included.
Write \(\alpha\) for the active parameters. For
\(v\in I\), let \(u_v\) be its cap axis on the central plane,
\(h_v=-\exp(\pi y_v u_v)\) to first order, and
\(t_{vw}=\epsilon_{vw}u_v\cdot u_w\) be a signed physical dot in
the marked seam frames. Put
\[
 B_I(u)_{vw}=
 \begin{cases}t_{vw}&vw\in E(I),\\0&\text{otherwise}.\end{cases}
\]
Signs can equivalently be moved to the axes since the graph is a forest.
All axis dots below refer to the common framed configuration, not to
independently chosen frames for different tests.

\begin{theorem}[All-angle cyclic Hall normal form]
\label{classic:thm:hall-action}
At every globally irreducible flat point in such a chart, the actual
finite harmonic Chern--Simons reduction is stably real-analytically
right equivalent to
\begin{equation}
 \mathrm{CS}_0+\frac{c}{2}y^{\mathsf T}B_I(u)y,
 \qquad |c|=\tfrac12 .
 \label{classic:eq:normal}
\end{equation}
Stability adds or removes only the finite gapped quadratic radial
pivots of this graph form. The equivalence fixes the entire central
section \(y=0\) pointwise in its literal angular and old active
parameters. It is jointly analytic through central active product
locks, including locally collinear locks, and through parallel as well
as perpendicular dormant angles. A choice of orientation fixes the
common sign convention for \(c\). The theorem does not identify all
off-shell extensions of the period marks under the normal change.
\end{theorem}

\subsection{Analytic representation scheme in peripheral marks}
\label{classic:sec:scheme}
We first describe the true flat ideal, before drawing any conclusion about
a potential. Use the graph-of-groups relations
\[
 [h_v,h_w]=1,\qquad A_v^2=h_v^{-1},\qquad
 {\rm Ad}_{A_v}P_v=P_v^{-1},
\]
with \(P_v\) the actual ordered product of transported neighboring fibres
(with their seam signs). The other cap root is
\(A_v^{-1}P_v^{-1}\). Local logarithms are taken after dividing off a
fixed central sign. All changes below use analytic units or invertible
analytic matrices and therefore hold in the analytic relation ideal,
not just on real points.

Write \(h_t=-\exp(p_tu_t)\) at an active centre \(t\). If an old active
neighbor of \(t\) is not parallel to \(u_t\), the corresponding
commutator kills \(p_t\) with a unit coefficient. At a remaining
\(t\), every incoming active axis is parallel to \(u_t\), and the
negative-cap equation at the base point implies \(P_t\) is central.
Thus each remaining column \(t\) touches at least two active
neighbors. An active \(s\) has \({\rm Ad}_{A_s}+1\) invertible; its
reversal equation kills \(\log(\pm P_s)\). Its longitudinal projection
has derivative a nonzero real entry on each incident surviving column
\(t\). This rectangular forest matrix has independent columns:
the graph induced by any nonempty supporting columns, with all their
neighboring rows, is a forest with every column degree at least two
and therefore has a row leaf. That row rules out its sole supported
column. Select a full-rank minor, including the earlier commutator
pivots. Every selected equation vanishes identically when \(p_T=0\).
Analytic division by that coordinate ideal and the invertible minor
therefore give the exact ideal \((p_T)\).

Once \(p_T=0\), all constraints at \(S\) are automatic. If \(P_t\)
is noncentral the scalar \(u_t\cdot\operatorname{Im}P_t=0\) solves
one angular coordinate. If \(P_t\) is central at the reference, use
the \emph{whole ordered product} as a local group parameter. To do
this simultaneously, root the product-incidence forest at old active
vertices. At each locked centre reserve one child active vertex;
its group element is solved from \(L h^\epsilon R=Q_t\) by actual
multiplication, including the dormant factors in their original
positions. Each reserved variable is used once. Hence all new
\(Q_t\) are independent analytic parameters. The locked constraint
is \(u_t\cdot\log(\pm Q_t)=0\); its two transverse product coordinates
are free. The lock is their joint zero, not an additional critical
factor. The result is an analytic product with the dormant critical
scheme and smooth passive \(\alpha\), even when several locks coincide.
A global effective rotation slice may be fixed at irreducibles.

For the all-negative dormant forest the nonparallel edge commutators
give \(q_e=p_vp_w\) in the ideal. The axial equation in the full
angular domain is
\[
 F_v=\sum_{w\sim v}b_{vw}(p,u)p_w,\qquad
 b_{vw}(0,u)=t_{vw}.
\]
Indeed, after removal of the fixed signs the product is an analytic
product of \(\exp(\pm p_wu_w)\). At a negative cap the longitudinal
projection of the reversal equation is identically a unit times its
axial logarithm. Multiply \(F_v\) by \(p_v\), remove all already known
\(q_e\), and peel the forest of edges parallel at the reference:
each leaf row has a single remaining coefficient with value \(\pm1\).
Thus all \(q_e\) belong to the ideal. Conversely the transverse
reversal equations vanish modulo \((q_e,F_v)\); there \(p_v\log P_v=0\),
and the transverse part is the pure half-turn equation. It follows
that \((q_e,F_v)\) is exactly the remaining ideal.

For use through parallel angles, write the row coefficients as
\(b_{vw}=t_{vw}+c_{vw}\), with \(c_{vw}\) in the ideal of \(v\)'s
neighboring periods, by radial integration of the ordered product
expansion. In \(\sum_w c_{wv}p_w\) each monomial is divisible by an
edge product \(p_wp_j\). Therefore
\(\widetilde t_{vw}=t_{vw}+c_{vw}+c_{wv}\) is a symmetric coefficient
whose radial rows agree with \(F_v\) modulo \((q_e)\). Differentiating
\(\sum_e\widetilde t_e q_e\) shows that its Jacobian is contained in
\((q_e,F_v)\). To reverse the inclusion use independent child-subtree
rotations on nonparallel reference edges, affecting only that edge's
literal dot, and add the rows \(p_v\partial_{p_v}\) on the parallel
forest. These rows generate all \(q_e\) by an invertible coefficient
minor: at zero periods the angular block is diagonal with nonzero
entries on the nonparallel edges, and the parallel subblock is the
same leaf-eliminable incidence matrix. The ordinary radial rows then
recover \(F_v\). Finally on parallel edges solve recursively
\(d_vd_w=\widetilde t_{vw}/t_{vw}\) with positive analytic units
\(d_v=1\) on \(p=0\). Set \(p'_v=d_vp_v\). On each nonparallel edge
use its independent relative subtree rotation to realize
\(\widetilde t_e/(d_vd_w)\) as the corrected literal dot. The
construction is invertible and fixes the central angular plane
pointwise. Thus the true dormant ideal in these ambient coordinates
is the Jacobian of
\[
 \frac12 p^{\mathsf T}B_I(u)p=\sum_e t_e p_vp_w .
\]
This is a graph-scheme statement only.

For clarity, the identification with the analytic Coulomb critical
scheme used here is the general holonomy correspondence on finite
local Artin rings \cite[\S\S5--6]{GoldmanMillson}, not a quadratic
formality theorem for this three-manifold. Fix an irreducible
effective slice on both sides. On a small Coulomb slice the
off-harmonic gradient equation has an invertible elliptic
linearization and is eliminated by the analytic implicit function
theorem (use Sobolev degree high enough for multiplication). The
remaining gradient represents the flat functor. The holonomy
comparison and the graph pivots identify these local quotient
analytic rings. In the ambient comparison, first split the radial
nondegenerate block of the graph Hessian at the reference. The
minimal rings then have the same tangent dimensions; lifts of
their analytic coordinate functions give inverse-function
isomorphisms of the ambient germs. They may be chosen to retain
the literal coordinate trace on the smooth common central section,
by correcting lifts by their trace there. Restore the finite
nondegenerate quadratic block on the Coulomb side. This gives a
stable coordinate comparison with the same Jacobian ideal and the
named central section, without claiming that equality of ideals
alone determines the Chern--Simons value away from it.

\subsection{Cyclic comparison of the actual radial Hessian}
\label{classic:sec:cyclic}
The extra input is a \emph{family} of symmetric radial differential
comparisons; comparing torsion numbers or signatures pointwise would
not suffice. First freeze the literal old active tuple \(\alpha_0\)
and let the dormant axes vary. Active pieces see only fixed central
dormant boundary holonomies on this plane, while each dormant piece
has a constant flat local system up to analytic rotation of its
cap axis. Seam changes of frame are their actual rotations. If the
active tuple itself is collinear, use a dormant nonparallel axis
for the residual azimuth; global irreducibility guarantees such an
axis. Thus freezing is a restriction on actual holonomies, not a
coordinate freeze that changes active holonomies when the dormant
axes move.

Use the twisted de Rham mapping cone of thickened pieces and collars.
There are no triple overlaps. The ordinary restriction map to this
cone is a chain map and induces the closed cohomology isomorphism
by the partition-of-unity proof of Mayer--Vietoris. Locally
analytic flat frames on a fixed cover can be patched with a smooth
partition of unity to identify nearby coefficient bundles; the
patched linear maps remain invertible near the base parameter and
polar decomposition preserves the inner product. Contract the
fixed-rank piece and seam complexes using analytic Hodge (or cellular)
contractions within this frozen family. If \(\delta\) is the
vertex-to-edge restriction difference and \(h\) is the block-diagonal
contraction, \(\delta h\delta=0\) because \(\delta\) kills the edge
summand. Thus the transferred differential is exactly the
cohomological restriction cone, with no higher-transfer terms.

Here is the precise radial part of its primitive elimination.
Use the common fibre class at each piece in its invariant space,
section periods \(c_{ve}\) with their primitive constraint
\(\sum_e J_{ve}^{*}c_{ve}=0\), and complementary one-sided area
lifts, where \(J_{ve}\) is the inclusion of the local invariant
space into the seam invariant space in the seam frame. Noncentral
torus complements are acyclic. At a central piece the
non-invariant part has independent restrictions on its central
exterior section planes, by the relative cohomology exact
sequence and the absence of invariant vectors there. Hence those
planes have unit elimination pivots. The swap identifies each
half-edge section with its neighboring fibre period. Substitution
in the section sum leaves the compressed signed adjacency.
In degree zero the inclusion difference on invariant vectors
and in degree two its adjoint are eliminated as injective and
surjective maps respectively, since the closed reference is
irreducible. They do not change the radial block. The connecting
cup formula is explicit in a collar: take a normal bump of
integral one and cup the associated section error with the
original fibre form. A radial error consequently pairs with its
own unit fibre period; a tangent with those dormant periods
identically zero pairs to zero. The ordinary determinant version
is also used in Theorem~\ref{thm:hall:torsion}.

If every \(T\) piece has no invariants, the radial block is
\(B_I\) plus the passive old active directions. At a locally
collinear lock its invariant line must instead be retained. Let
\(T_1\) consist of those lines in the frozen family. In the
\(S,T_1,I\) order the radial block is exactly
\begin{equation}
 \begin{pmatrix}0&U^{\mathsf T}&0\\
                 U&C&V\\
                 0&V^{\mathsf T}&B_I
 \end{pmatrix}.
 \label{classic:eq:activeblock}
\end{equation}
The entries of \(U\) on the incidence forest are signed unit
axis pairings; every \(T_1\) row has at least two old active
neighbors. The leaf proof above again shows \(U^{\mathsf T}\)
is injective. Split \(S=S'\oplus\ker U\) so the restriction
\(U'\) to \(S'\) is square invertible. The inverse of the
\(S'\oplus T_1\) block is
\[
 \begin{pmatrix}
 -(U')^{-1}C((U')^{\mathsf T})^{-1}&(U')^{-1}\\
 ((U')^{\mathsf T})^{-1}&0
 \end{pmatrix}.
\]
In particular it has \emph{zero hub--hub block}. The dormant
coupling reaches only that block, so the Schur correction is
identically zero and the retained graph is \(B_I\) itself.
A dormant tangent may acquire an \(S'\) correction, but its
cup pairing with a dormant connecting error is still the unit
dormant period. This is why local collinear locks do not induce
a new dormant spring or rescale its radial units.

On the true closed Hodge side, an analytic sufficiently small
positive gap for the twisted Laplacian around the irreducible
reference isolates degrees one and two, with star identifying
the two low spaces. The analytic Hodge projection away from
exact forms uses the nonzero scalar gap. The tangent to the
central family gives analytic harmonic modes in the kernel of
the corresponding coexact curl; split these angular modes,
using the star-dual split in degree two. Similarly split the
already nondegenerate radial block of \(B_I\) at the reference
by ordinary analytic symmetric Gaussian elimination. The closed
low complex and the model cone now give analytic chain maps
\(T_1',Q_2'\) on their remaining radial blocks:
\[
       B\,T_1'=Q_2'D ,
\]
where \(B\) denotes the graph Schur block and \(D\) the low
symmetric differential in physical units after the fixed
conversion described below. The maps on radial blocks are
invertible at the reference and hence nearby. At the base
point the cone map is a cup isomorphism; in the splits just
described the degree-one map sends angular tangents to angular
tangents, which pair to zero with radial degree-two connecting
errors. The induced quotient pairing of the radial tangents
and radial connecting errors is the unit one by the primitive
torus computation. Thus
\((T_1')^{\mathsf T}Q_2'=1\) at the reference. No pointwise
minimal basis is chosen as a function of the dormant rank.

To compare with the \emph{nonlinear} harmonic reduction's Hessian,
project from the static to the moving Coulomb slice along exact
directions; on a flat connection the Hessian kills those
directions. The static transverse high complement is a small
analytic graph \(J\) over the moving high spectral complement.
Writing the moving forms as \(D_L,D_H\), its static Schur
Hessian is
\[
 D_L-D_LJ(D_H+J^{\mathsf T}D_LJ)^{-1}J^{\mathsf T}D_L.
\]
The gapped inverse is analytic in the elliptic Sobolev spaces
and the correction acts through a finite-dimensional space.
This form is \(A D_L\) with \(A\) analytic, equal to the
identity at the reference and symmetric after multiplication
by \(D_L\). Accordingly the analytic square root gives an
actual congruence. The same square-root argument applies to
the cone comparison: \(P=(T_1')^{\mathsf T}Q_2'\) is near
one and \(PD=(T_1')^{\mathsf T}BT_1'\) is symmetric, so its
analytic principal root obeys
\(P^{1/2}D=D(P^{1/2})^{\mathsf T}\), even at rank changes.
The Schur complement comparison extends back to the
stabilized radial block by using the explicit nondegenerate
pivot. Finally the primitive fibre period in \(y\) is
\(\pi\,dy\), its degree-two partner is cup-dual, and the CS
Hessian is \(\pm (2\pi^2)^{-1}\) times the coexact curl
pairing, obtained by differentiating
\((8\pi^2)^{-1}\int\operatorname{Tr}(A\,dA+2A^3/3)\).
The symmetric Hessian in signed Wilson units is thus
\(\pm\tfrac12 B_I\), independent of any WRT coefficient.

It remains to extend the frozen comparison analytically in
\(\alpha\), as the cut-piece cohomology need not remain
constant. Write \(F(y,u,\alpha)\) for the actual stabilized
germ in the common Jacobian chart, with its constant flat value
subtracted and \(F_2=\tfrac12 y^{\mathsf T}D(\alpha,u)y\).
Let \(G=\tfrac{c}{2}y^{\mathsf T}B_I(u)y\). The shared
Jacobian ideal is generated in radial degrees one and two
by the linear rows \(B_Iy\) and all edge products. At
\(\alpha_0\) use the analytic radial congruence just obtained.
It preserves the linear ideal by the degree-one equality.
It also preserves the degree-two edge module modulo the
linear ideal: \(\operatorname{Jac}F_2\subseteq
\operatorname{Jac}F\) by radial homogeneity in the common
graded ideal. The induced invertible symmetric-square map
on quadratic polynomials therefore sends that module into
itself. Its inverse is a polynomial in the map, by
Cayley--Hamilton with unit constant term, so the inclusion
reverses. This argument does not require the edge module to
be locally free. Apply this one \(u\)-dependent coordinate
change to all active parameters. Now
\(D(\alpha,u)=A(\alpha,u)cB_I(u)\) for analytic \(A\), by
the linear ideal equality, and \(D(\alpha_0,u)=cB_I(u)\)
for every \(u\) in the chart. A forest with a perfect matching
has determinant \((-1)^d\prod_{e\in M}t_e^2\) (only the
matched two-cycles contribute). It is invertible at a dense
set of literal angle configurations, including in a chart
through a perpendicular or parallel reference. Hence
\(A(\alpha_0,u)=1\) identically and its principal square root
gives the analytic congruence jointly in \((\alpha,u)\).
This avoids passing a cut-piece contraction through a jump.

Lastly the radial-quadratic truncation has the same full
Jacobian by the preceding degree-two argument. Use the
straight homotopy \(F_s=G+s(F-G)\) after the congruence.
The gradient differences express in the homogeneous graph
generators with coefficient matrix whose value at \(y=0\)
has only the radial-row/angular-column block possibly
nonzero. Indeed the radial differences start at degree two,
angular differences at degree three, and each edge product
is a combination of angular rows and \(y\) times radial
rows, as in the nonparallel/parallel pivots above. Thus
the coefficient matrix for the gradients of \(F_s\) is
uniformly invertible near \(y=0\), since its constant
correction is strictly block triangular. Radial Euler
integration of \(F-G\) supplies coefficients of order
\(O(y^2)\) on radial gradient rows and \(O(y)\) on angular
gradient rows. Applying the inverse preserves these orders.
The analytic Moser vector field solving
\(dF_s(X_s)=-(F-G)\) therefore has \(X_y=O(y^2)\) and
\(X_{u,\alpha}=O(y)\). Its local flow fixes the entire
central section; this proves Theorem~\ref{classic:thm:hall-action}.
The preliminary radial congruence need not be the identity
on all off-shell period directions.

The ideal elimination at old active centres also works
when some of them are positive. Use the three-vector
period at a positive vertex; each genuinely noncentral
neighbor kills the two transverse components, and any
remaining axial column has at least two old active
neighbors (its product must be central). The longitudinal
forest minor again eliminates them all, and the locked
positive product target is fixed in all three components.
This is only a separation of the framed critical scheme.
A positive \emph{dormant} component has vector periods and
is not represented by \eqref{classic:eq:normal}.

\subsection{Absolute density and what the leading angular law measures}
\label{classic:sec:leading}
Put \(a=|T|\), \(b=|S|\), \(2d=|I|\), \(N=a+b+2d\)
and \(\nu=a+d\) on the intrinsic Hall family. Let
\[
 p_G=\frac{3N-2\nu-3}{2},\qquad
 C_G=2^{(3N+1)/2-2b-3d}\pi^{1-N}.
\]
Write \(dM_*\) for the forced radial/ordered-product law of
\eqref{eq:hall:weight}, and \(dP\) for its accompanying common
normalized angular law. In particular every dormant seam dot
has conditional marginal \(dt/2\) before test weights are inserted:
a successive relative rotation along a tree produces fresh spherical
probability independently of the previous seam dots. The other
angles are \emph{conditioned jointly}, not recreated separately
for each insertion. Let \(dP_\perp\) denote this law with the
matched edge dots set to zero by conditional probability.

Contract the unique dormant matching and write an unmatched
edge \(a_i b_j\) as \(i\to j\), for a fixed bipartition within
each component. The order polytope \(\Omega\) has one
coordinate per matched pair, inequalities \(u_i\le u_j\)
along these arrows, and \(0\le u_i\le1\).
In the classical Fourier proof one can use
\(u_i=1+\log|y_{a_i}|/\log r\); the use of
\(-\log|y_{a_i}|/\log r\) instead reverses the inequalities
and has identical volume.

\begin{theorem}[The finite canonically normalized Hall coefficient]
\label{classic:thm:hall-stationary}
On the irreducible all-negative Hall charts of
Theorem~\ref{classic:thm:hall-action}, the leading angular
density of the finite CS residual model, after removing the
phase \((-1)^{r(N-1)}\) and \(r^{p_G}(\log r)^d\), is
\begin{equation}
 C_G\operatorname{Vol}(\Omega)\,dM_*\,dP_\perp .
 \label{classic:eq:hallleading}
\end{equation}
This is valid for fixed smooth tests on compact chart
localizations; it patches on compact Hall interiors. If
\(d>0\), a fixed smooth test whose trace on the indicated
matched angular support vanishes is \(O(r^{p_G}(\log r)^{d-1})\)
on a fixed compact localization. Exhausting the resolved
leading measure in the forced variables gives the entire
vacuum coefficient in Theorem~\ref{thm:hall:main}.
Equality of those totals does not assign every other classical
stratum a zero contribution.
\end{theorem}

\begin{proof}
The invariant Gaussian factor for \(h\) retained irreducible
zero modes, in cohomology frames of the stated metric, is
\(\tfrac12(r/(2\pi^2))^{h/2}\sqrt{T_{\rm adj}}\); the factor
\(1/2\) divides by the ineffective central stabilizer, not
by a continuous rotation volume a second time.
This is the quadratic integration and framing normalization
of \cite[eqs.~(3.3), (3.6)--(3.7), (3.14), (4.5)]{Adams}.
Use the analytic spectral gap before separating small eigenvalues:
its complement determinant is smooth as flat parameters vary.
If \(L\) is the Hessian of the finite reduced action on the
normal quotient, the corresponding small coexact curl product
has square root
\((2\pi^2)^{\dim L/2}\sqrt{|\det L|}/\sqrt{\det G_N}\),
where \(G_N\) is the normal quotient metric. One way to see
the precise relative statement near degeneracy is the
triangular Schur graph of the high block, followed by the
invertible low spectral projection; it intertwines the
bilinear forms up to factors tending to identity, rather
than by an additive estimate of individual eigenvalues.
The coordinate volume splits off exactly
\(\sqrt{\det G_N}\) by the Schur identity. Thus the normal
metric cancels out of the determinant-line restriction.
This gives the usual stationary determinant law here
\emph{also} as one approaches the rank-changing points
through the regular chart.

The independent regular coarea computation in
Theorem~\ref{thm:hall:torsion} gives the coefficient with
\(r^{p_G}\) removed,
\[
 D_G\,\frac{dM_*\,dP}{\prod_{e\in M}|t_e|},
 \qquad
 D_G=(2\pi^2)^{-p_G}\,2^{a+N-1}\pi^{2b+2d-2}.
\]
The stabilized graph's radial Gaussian at generic angles has
factor \(2^d/(r^d\prod_M |t_e|)\) because \(|c|=1/2\),
\(|\det B_I|^{1/2}=\prod_M|t_e|\), and its signature is zero.
Hence the smooth stabilized amplitude at zero radii is
\(2^{-d}D_G\,dM_*\,dP\), with common \(r^{p_G+d}\).
It agrees first on an open dense angular set by the
determinant comparison and extends therefrom by smoothness.
For all negative fibres/centres the CS action is
\((N-1)/2\bmod1\) and the relative adjoint eta phase is zero
by the phase atlas. This also removes any putative extra
signature from the comparison Gaussian.

For the leading singular integration write
\(x_i=y_{a_i}\), \(z_j=y_{b_j}\).
First Fourier integrate the \(z_j\)'s of a smooth compact
amplitude at the frequencies
\(r c\sum_i t_{ij}x_i\). Repeated integrations by parts
give arbitrary product decay in these frequencies. Integrate
one dominant dot per \(b_j\); these dots are distinct and
their product marginal is bounded. The resulting bound is
\(\prod_j C/(1+r\max_{i\sim j}|x_i|)\).
Each violating order ratio \(|x_i/x_j|>1\) then costs
the integrable reciprocal ratio, thus at least one logarithm
relative to \(d\) independent \(d|x_i|/|x_i|\) scales.
At the short radial end \(|x_j|\lesssim r^{-1}\) the
same bound gives one fewer log, and the finite upper end
has the same property. On the ordered chamber set
\(w_j=rx_j z_j\). Extending the compact radial domain in
these Fourier coordinates and replacing the radial amplitude
by its central trace have one-log error: the bounds
are integrable with relative factors \(|x_i|\) and
\((r|x_j|)^{-1}\), or the rapid Fourier moments of
\(|w_j|/(r|x_j|)\). Writing the matched-frequency relation
as
\[
 t_{jj}+\sum_{i\ne j}(x_i/x_j)t_{ij}=0
\]
shows that its matched-dot centres shift by
\(O(\sum |x_i/x_j|)\), again integrable with one logarithm
lost. One may first discard a bounded-width log strip near a
tie to keep these centres in the interior of every matched
dot interval. Fourier inversion there gives \(|c|^{-1}\)
per matched dot and the product density \(dt/2\) supplies
\(1/2\) per dot. This uses ordinary smooth conditional
densities near the \emph{matched} dot zero only. Nonmatched
angle poles can be handled by using spherical frames and
bounded smooth tests throughout, without a coordinate
singularity at a pole. The signed \(x_i\) ranges supply
\(2^d\); the ordered absolute log region has volume
\((\log r)^d\operatorname{Vol}(\Omega)+O((\log r)^{d-1})\).
Multiplying these factors leaves \(D_G=C_G\) and the
joint law \eqref{classic:eq:hallleading}, with the stated
loss if its trace vanishes. For \(d=0\) this step is
just regular evaluation. Relative exhaustion of the
finite forced mass and its equality with the exact
vacuum coefficient are the independent quantum estimates
in the Hall theorem; the two limits are not interchanged.
\end{proof}

The smooth product resolution makes a negative active lock
a real two-vector zero. Its mass on a compact fixed-rank
resolved chart is \(O(\epsilon^2)\) when that vector lies
in a ball of radius \(\epsilon\); a projection to active
radii can instead have a cap-convolution singularity.
One should use its independently resolved Haar estimate
rather than infer bounded projected density from the smooth
chart. When \(d>0\), a globally collinear all-negative
Hall reference has \(|t_e|=1\) on every dormant edge;
the graph block is invertible there and no dormant
logarithm arises in its fixed compact residual gauge model.
For \(d=0\) a collinear orbit has a circle stabilizer rather
than a smooth ordinary quotient. In the fixed slice use
Haar division before quotient: a generator of orbit speed
\(v\) and its cup dual contribute a torsion square-root
factor \(|v|\) on the nonfixed locus, the orbit-volume
Jacobian of this division. The effective orbit period is
half the circle's full period. This local pushforward
adds no atom on fixed orbits. These observations are not
uniform estimates near a change in the noncentral set.

\subsection{Tied scales carry a different angular support}
\label{classic:sec:overlap}
There is a direct subleading test of the same actual
normal phase. Choose any subset \(E'\) of unmatched arrows
and insert compact smooth cutoffs on the signed ratios
\(\rho_e=x_i/x_j\in\mathbb R^*\) for \(e=i\to j\) in \(E'\).
The constrained pairs form a forest. Let \(p=d-|E'|\) be
the number of its components, retaining isolated matched
pairs, and let \(\Omega_{E'}\) be the order polytope of the
roots under the unselected arrows. In any chart with smooth
density in the matched dots on the selected support (in
particular compact interior angular traces), the same
finite model has order \(r^{p_G}(\log r)^p\); its
coefficient is obtained by the replacements
\[
 \begin{gathered}
  2^d\operatorname{Vol}(\Omega)
      \longmapsto 2^p\operatorname{Vol}(\Omega_{E'}),\\
  \delta(t_{jj})\longmapsto
  \delta\!\left(t_{jj}+
       \sum_{\substack{i\to j\\(i,j)\in E'}}\rho_{ij}t_{ij}\right),
  \qquad
  \prod_{e\in E'}\chi_e(\rho_e)\frac{d\rho_e}{|\rho_e|}.
 \end{gathered}
\]
Here the densities and Fourier modulus per matched dot
are the same as before, a delta outside the angular
cutoff contributes zero, and the error has one fewer
root logarithm on fixed compact supports. Indeed, change
the absolute log coordinates to one root per selected
component and one edge difference per selected edge.
The incidence minor has determinant of modulus one;
the selected ratios fix descendant signs so only \(2^p\)
root signs remain. Unselected ratios have the same
dominant-bound and Fourier-shift estimates as above.
The selected coherent frequencies cannot be put equal
to zero in the integrand. Fourier inversion gives the
displayed deltas with unit matched-angular derivative.

If all arrows in one dormant connected component are
selected, the delta relations exhibit a genuine flat
curve: set all \(z\)'s to zero, let the \(x\)'s be a
common scalar times the vector of the selected signed
ratios, and keep the indicated \(u\) fixed. Every edge
period product vanishes and its only nontrivial radial
row is the displayed balance. The inverse analytic
chart therefore gives actual flat holonomies with this
physical first-order direction. A single selected arrow
in \(P_4\) can already have
\(t_{11}=0,\ t_{22}=-\rho t_{12}\ne0\), away from the
principal perpendicular-pair support. This is not a
positive splitting of the full unmarked next
coefficient; the normal off-shell ratios are not
automatically the same tests as the literal finite
quantum ratio marks.

\subsection{The absolute two-vertex and four-vertex tests}
\label{classic:sec:benchmarks}
These computations compare independent quadratic normalization and
exact finite level, so in particular check the volume of the effective
rotation group and the finite centre factor. The all-negative
\(P_4\) test also checks why an apparently singular intermediate
angle does not create an extra leading logarithm.

\paragraph{Two sites: the flat \(H\) model.}
The two-site \(H\) is the didicosm. Its degree-four rectangular torus
cover has lengths \(L_1,L_2,L_3\); the three nontrivial deck motions
have linear half-turns about the coordinate axes and half-period
translation on their fixed axis. These data, and the orthogonal
quaternionic flat lift, follow from the two glide presentations
in \cite[\S2]{P49}. Put \(V=L_1L_2L_3/4\).
The twist is the three-dimensional linear representation of the
deck quotient \(V_4\). Its harmonic one-forms are
\(X\mathbf i\,dx_1+Y\mathbf j\,dx_2+Z\mathbf k\,dx_3\) and
there is no invariant vector. These parallel forms close under the
quadratic curvature followed by Hodge star, so the complement
equation is solved identically by zero. The residual CS magnitude
is \(V|XYZ|/\pi^2\), directly by the normalized trace of the cubic.

The torsion can be computed without a seam convention. In the
weighted heat trace
\[
 \sum_q(-1)^q q\operatorname{Tr}(e^{-s\Delta_q}),
\]
the identity deck insertion has zero weighted form character.
A half-turn has form character \(-4\), and
adjoint character \(-1\); deck averaging gives \(+1\).
Its restricted Fourier sum is
\(\sum_{n\in\mathbb Z}(-1)^n e^{-s(2\pi n/L_i)^2}\).
Omitting its zero term gives in zeta variables
\[
 \zeta_i(z)=2(L_i/(2\pi))^{2z}
       (2^{1-2z}-1)\zeta_{\rm R}(2z),\qquad
 \zeta_i'(0)=4\log2-2\log L_i .
\]
Therefore \(T_{\rm adj}=\exp(\frac12\sum_i\zeta_i'(0))=16/V\).
The coexact curl has equal opposite helicity multiplicities in
both twisted and untwisted deck traces: on any nonidentity fixed
wavevector the half-turn acts with the same sign on both
transverse helicity lines. The scalar/exact parts also pair,
so the relative eta is zero. The harmonic coordinate metric
has volume \(V^{3/2}\,dX\,dY\,dZ\). The normalized zero-mode
amplitude is thus \(r^{3/2}V/(\sqrt2\pi^3)\,dX\,dY\,dZ\).
The cubic with \(\mathrm{CS}_0=\frac12\) gives
\((-1)^r\sqrt{2r}\log r/\pi\). Here and below a split
cubic \(cXYZ\), with a smooth cutoff around zero, has
coefficient \(2/(r|c|)\log r\) times the amplitude at zero:
Fourier transform one pair, integrate the remaining
signed \(dX/|X|\) between \(r^{-1}\) and a fixed radius,
and bound the two endpoint errors by the same rapid
Fourier estimate as in the Hall proof. This agrees
with the independent finite \(H\) computation in \cite{P48,P49}.

\begin{theorem}[Two normalized singular \(P_4\) models]
\label{classic:thm:p4}
For the all-negative central \(P_4\), let \(s_1,s_2\) be its
two outer cap-axis dots and \(t\) its middle dot. At
\(s_1=s_2=0\) and \(0<|t|<1\), the physical reduced
quartic magnitude in outer signed Wilson units is
\[
 \frac{|s_1s_2 y_1y_2|}{2|t|},\qquad
 \sqrt{T_{\rm sing}}=
 \frac{2\pi^4}{|t|}\,ds_1\,ds_2\,dy_1\,dy_2\,dt\,d\gamma .
\]
Here \(d\gamma\) denotes two passive azimuth probabilities
and \(\sqrt{T_{\rm sing}}\) is the coarse half-torsion.
After canonical gauge/framing normalization the leading
central model and the independent exact finite middle-dot
orbital law each give
\[
 \frac{(-1)^r\sqrt2}{4\pi^3}\,
      r^{5/2}(\log r)^2\,dt
\]
after passive angular integration. Hence the closed coefficient
is \((-1)^r\sqrt2/(2\pi^3)\) at \(r^{5/2}(\log r)^2\)
in \(Z_r\). In the literal disjoint trace box with signs
\((-,+,+,-)\), there is a phase-zero singular cubic at
perpendicular negative endpoint caps. Its total normalized
coefficient is \(+\sqrt2/(16\pi)\) at \(r^{1/2}\log r\),
matching the exact finite trace-box contraction.
\end{theorem}

\begin{proof}
For the all-negative corner, eliminate the two central
radial variables in the exact graph action:
the quadratic graph expression before its common coefficient
\(c\) is \(tX_1X_2+s_1X_1y_1+s_2X_2y_2\).
Its stationary value is \(-s_1s_2y_1y_2/t\), with retained
signed outer period linear units on the singular stratum.
For an independent absolute check, an all-central negative
Seifert piece contributes adjoint torsion 4 in primitive
peripheral bases. To see this, remove its genuine caps.
The base--fibre product cochain complex cancels against its
one-degree shift in the paired section/fibre bases. On the
plane perpendicular to the cap axis both cap cores act by
\(-1\), giving determinant 4 on each cap solid torus; their
tori are acyclic on that plane. On the invariant line the
two degree-one cap restrictions have determinants 2 and
2 with unit section pivots. The degree-zero and degree-two
pivots have modulus one; the opposite slopes cancel in
the surviving section relation. Thus the local product is
\(4^2/2^2=4\). The trivial positive central block, used
below, has instead three line copies and torsion \(4^{-3}\).

The same relative cone computation as in the Hall proof
then leaves two angular maps \(U,U^*\) and the dot
adjacency in degree \(1\to2\); the four piece contributions
have product half-factor \(2^4\). The radial pivot scalar
is \(|t|\) in half-torsion, and the retained two period
units give \(\pi^2\). For the angular determinant, lift
axis variations to \(\delta u_v=2a_v\times u_v\). Edge
differences plus the root generator are a triangular
unimodular change on the tree. Removing axial columns
and the global rotation leaves
\[
  \frac{\prod_{v=1}^4 dA(u_v)}{4^4\pi^2}
       =\frac{\pi^2}{8}\,ds_1\,ds_2\,dt\,d\gamma
\]
as a local quotient density, the right side using three
uniform edge dots. This yields \(2\pi^4/|t|\) as stated.
There is no second quotient by \(SO(3)\).
The clean outer-active branch is an independent
determinant check: per signed quadrant its coarse
density is
\(\pi^2\cot(\pi |y_1|/2)\cot(\pi |y_2|/2)
\,d|y_1|\,d|y_2|\,d\operatorname{Vol}_{SO(3)}\),
where \(d\operatorname{Vol}_{SO(3)}=(\pi^2/2)dt\,d\gamma\).
Its residue is \(2\pi^2/|y_1y_2|\). Dividing the singular
coarse density by the small curl product
\(2\pi^2|\kappa y_1y_2|\), with \(\kappa\) the physical
quartic coefficient, gives this residue exactly for
\(|\kappa|=1/(2|t|)\). The small eigenvalues have opposite
signs; the phase is \((-1)^r\) under
Section~\ref{phase:classical}.

For the finite quartic \(s_1s_2y_1y_2\), Fourier integrate
\(s_1,s_2\) at fixed signed \(y_1,y_2\). The four signed
quadrants each have logarithmic area \(1/2\) after the
gate \(r|y_1y_2|\gtrsim1\), yielding coefficient
\(2/(r|\kappa|)\,(\log r)^2\). The omitted gate strip and
the lower radial ends are controlled by the same Fourier
moments and lose one log. Combining the \(h=7\) Gaussian
normalization with the independent singular density leaves
the given constant multiple of \(dt\).

There is a global exact angular check without uniform
classical estimates at \(t=0,\pm1\). Use a smooth even
middle orbital test \(\eta(t)\) and its kernel
\(K_\eta(\xi)=\frac12\int_{-1}^1\eta(t)e^{i\xi t}dt\).
In the finite reflected deficit sum eliminate both
outer leaves by the odd sine identity
\[
 \sum_{\ell<r,\ \ell\ {\rm odd}}
     \cot\!\bigl(\tfrac{\pi\ell}{2r}\bigr)
                    \sin\!\bigl(\tfrac{\pi n\ell}{r}\bigr)
     =\tfrac r2,\qquad 0<n<r,\quad n\ {\rm odd}.
\]
For instance the identity follows by subtracting adjacent
odd frequencies using sine/cosine addition on the odd
grid and evaluating the first frequency. Using the
prefactor and vertex factors of \eqref{rad:sum}, the
remaining exact expression is
\[
 (-1)^r\frac{4\sqrt2}{\pi^3}r^{5/2}
 \sum_{\substack{n,m<r\\ n,m\ {\rm odd}}}
 \frac{F_2(n/r)F_2(m/r)}{nm}\,
 K_\eta(\pi nm/r),\qquad
 F_2(x)=\left(\frac{\pi x/2}{\sin(\pi x/2)}\right)^2 .
\]
Indeed the leaf identity in sinc units is \(r/(4n)\);
the middle vertex has precisely the bounded endpoint
factor \(F_2\). The normalized odd harmonic measures
converge to \(\tfrac12\,du\) with \(u=\log n/\log r\).
Away from the endpoint strips and \(u+v=1\), the
bounded kernel tends uniformly to \(K_\eta(0)\) on
one triangle and 0 on the other. The strips have small
normalized mass and \(F_2\) tends to one on interior
logarithmic scales. The triangle's weight is \(1/8\).
This gives the exact claimed density for even orbital
tests, including exhaustions of the open generic
angular set. The check does not substitute an orbital
integral variable for a genuine finite Wilson operator.

For the mixed sign box, fix one root sign at each positive
site. Use signed radian endpoint periods \(p,q\), axes
\(u,v\), and \(t=u\cdot v\). The positive centres' vector
periods are eliminated by reversing their ordered
products: their caps have \({\rm Ad}+1\) invertible,
forcing the two middle periods to be \(-qv,-pu\).
The edge commutators then give \(pq=0\), with a unit
from a component of \(u\times v\) at perpendicular
axes. The two negative caps' axial equations give
\(pt=qt=0\); the transverse equations hold modulo these
generators. Thus the exact critical scheme is
\(\operatorname{Jac}(pqt)\). The genuine CS germ has no
quadratic terms in these three retained coordinates.
Only \(pqt\) can occur in its cubic with these gradient
generators, and its coefficient cannot vanish.
The analytic straight-homotopy argument for this
graded ideal gives the exact cubic with a
tangent-identity change (equivalently the normal-form
argument of \cite{P49B}); there is no change in the
physical period units. \(\mathrm{CS}_0=0\) since the
two positive product bits are both one, and the
relative rho vanishes by \eqref{phase:rho}.

The singular half-torsion at \(p=q=t=0\) per root pair
is \(dp\,dq\,dt/32\). In fact the four blocks now give
product half-factor \(2^2(1/8)^2=1/16\).
The degree-one compressed adjacency has half-block
\((x,b)\mapsto(u\cdot x,x+bv)\); its nonzero singular
values are \(\sqrt2,\sqrt2,1\).
The two retained primitive radial tangent vectors
have norm \(\sqrt2\) and are orthogonal, cancelling
the determinant divisor 2. The two axis areas
divided by \(4^2\pi^2\) have total weight one,
with normalized dot density \(dt/2\). This verifies
the stated singular factor. On the adjacent reflected
smooth arm \(p\ne0,\ q=t=0\), the coarse half-torsion
per root pair tends to \(dp/(32|p|)\).
The compact-regular formula is
\((\pi/64)\cot(\pi x/2)\csc^2(\pi(1-x)/2)\,dx\)
with \(x=|p|/\pi\). Its primitive block computation
has two ordinary noncentral factors
\(\frac12(4\cos^2(\pi x_j/2))
 (4\sin^2(\pi x_j))^{-d_j/2}\),
the positive centre factor \(1/8\) and the negative
end factor 2; the invariant radial path block has
the common unit period in radians and all
transverse section/area pivots have determinant
one in coarea bases. Substituting \(d_j=1,2\) and
\(x_j=x,1-x\) gives the displayed formula.
The small-curl determinant comparison now yields
\(|\kappa|=1/(2\pi^2)\) in \(pqt\) radian units.
With the \(h=3\) canonical factor and the split
cubic residue, the coefficient is
\(1/(32\sqrt2\pi)\) per positive-root pair. The four
such pairs are distinct. Ordered perpendicular
negative axes form a single effective rotation
orbit; the two signed period variables already
include both adjacent Weyl directions. The total
is \(+\sqrt2/(16\pi)\).

For comparison the exact finite trace box can be
reduced by the same two leaf sums with cutoffs.
After reflecting the two small positive-centre
labels \(a,b\), the middle sine changes by
\((-1)^r\), cancelling the external parity.
The remaining form is
\[
 \frac{4\sqrt2}{\pi^4}r^{-1/2}
 \sum_{a,b\equiv r-1\, (2)} C_2(a/r)C_3(b/r)
   \sin(\pi ab/r)\,Q_1^{(r)}(a)Q_4^{(r)}(b),
\]
where the positive-centre profiles are smooth with
\(C_j(0)=\pi^2/4\), and cut leaves give
\(Q_i^{(r)}(a)=1+O(1/a)\) on the sufficiently small
positive arguments. The latter estimate follows by
summation by parts in the smooth difference from
the uncut cotangent identity. For \(a\) in this range,
two geometric summations by parts in \(b\) give
\[
 \sum_{b\equiv r-1\, (2)} C(b/r)\sin(\pi ab/r)
    =\frac{r\,C(0)}{2\pi a}+O_C(r/a^2+1):
\]
the first endpoint term is the imaginary part of
\(C(s/r)e^{is\pi a/r}/(1-e^{2i\pi a/r})\),
\(s=1\) or 2, and the second difference has total
variation \(O(r^{-1})\). The single leaf errors cost
\(O(r\sum a^{-2})\), the double error \(O((\log r)^2)\).
Now the half-grid harmonic sum gives \(\frac12\log r+O(1)\)
on either parity. The result is exactly
\(+\sqrt2\,r^{1/2}\log r/(16\pi)+O(r^{1/2})\).
It is a literal trace box, not a finite projector onto
a single character, and two independently integrated
smooth endpoint arms must not each be credited with
the whole logarithm.
\end{proof}
